% ============================================================
%  IV. WHOLE-BODY ACTIVE PERCEPTION
% ============================================================
\section{Whole-Body Active Perception}
\label{sec:active_perception}

The core contribution of \ac{TERESA} is a whole-body active perception
framework: the robot continuously adapts its posture---body height, pitch,
and yaw---to maximise the quality of its perceptual input. This section
describes the three mechanisms that realise this principle:
confidence-gated active search (Section~\ref{sec:confidence_search}),
anticipatory body scanning with \ac{WBC} (Section~\ref{sec:anticipatory_scan}),
and perception-quality-aware velocity control
(Section~\ref{sec:quality_velocity}).

\subsection{Confidence-Gated Active Search}
\label{sec:confidence_search}

Before the robot can approach the patient, it must first localise them
reliably. In an unstructured emergency scenario, the patient's position is
unknown a priori, and the detection quality of the onboard Orbbec \ac{RGBD}
camera depends strongly on the robot's body posture: a higher camera provides
a wider field of view but reduces resolution at the ground plane; a forward
pitch directs the camera toward the expected patient location but narrows the
horizon. \ac{TERESA} treats this coupling as an opportunity rather than a
limitation: the robot actively searches over body postures to find the
configuration that maximises detection confidence.

\subsubsection{Posture grid search}
The robot enters a SEARCHING phase in which it explores a discrete grid of
body postures. At each grid point, Spot assumes a lowered body height and a
target pitch and yaw, pausing for a fixed interval to allow the Orbbec
perception pipeline to process several frames. The postures are generated by
varying yaw (orientation about the vertical axis) and pitch (forward
inclination) over a set of pre-defined discrete angles, while keeping the
body lowered to bring the camera closer to the ground plane where a lying
patient is expected.

\subsubsection{Confidence lock}
The Orbbec pipeline runs \ac{YOLO}-based skeleton detection and classifies
the patient's posture, producing a confidence score $c(\boldsymbol{\xi},
h, \varphi) \in [0, 1]$ that quantifies the reliability of the detection.
When this confidence exceeds a high threshold, the robot immediately stops
the posture grid and freezes in place. It then collects a short sequence of
patient position estimates, transforms them from the camera frame to the
world-fixed odometry frame $\mathcal{F}_W$ via the \ac{TF} tree, and computes
their mean as the initial approach target. This confidence-gated lock ensures
that the robot commits to an approach trajectory only once the patient has
been localised with high certainty, avoiding premature navigation toward
noisy or incorrect detections.

\subsubsection{Robustness to confidence loss}
If the detection confidence drops below the lock threshold during target
sampling---for instance, due to a temporary occlusion or a change in
lighting---the robot abandons the lock and resumes the posture grid from the
current configuration. This closed-loop behaviour prevents the system from
committing to a stale or degraded target.

Once the target is locked and sampled, the robot straightens its posture
(returning to nominal body height and zero pitch), and the mission transitions
to the approach phase. During approach, a second perception pipeline on the
\ac{EE}-mounted RealSense camera takes over for fine-grained patient analysis.

\subsection{Anticipatory Body Scanning with Whole-Body Control}
\label{sec:anticipatory_scan}

In traditional mobile manipulation pipelines, the robot first navigates to a
target position, stops, and only then begins visual scanning---a sequential
approach that introduces significant idle time. \ac{TERESA} eliminates this
idle period by executing the body scan concurrently with the approach
navigation. This anticipatory scanning is enabled by a decoupled architecture
in which the Spot base navigates toward the patient while the Z1 arm
simultaneously performs a multi-view scan of the patient's torso.

\subsubsection{Decoupled whole-body coordination}
During the approach phase, two independent controllers run in parallel. The
Spot base is driven by a P-controller that computes forward velocity
$v_x = k_p^{\mathrm{lin}}\,\Delta d$ and yaw rate
$\omega_z = k_p^{\mathrm{ang}}\,\Delta\theta$ from the distance and angle
error between the current body pose and the target in $\mathcal{F}_W$. The
arm is controlled by a \ac{WBC} \ac{QP} that computes joint velocities
$\dot{\mathbf{q}}$ via a damped pseudo-inverse of the arm Jacobian,
generating an EE orientation that points the RealSense camera toward the
patient's torso. The \ac{WBC} is used here not merely to optimise
kinematics, but to serve the active perception objective: it continuously
reorients the sensor to maximise the quality of the anatomical data being
acquired.

\subsubsection{Multi-view ARC scan}
The arm executes an ARC\_GRID scan: a sequence of pre-defined poses that
sweep the RealSense camera across the patient's torso from multiple
viewpoints. Each pose combines a wrist angle (to tilt the camera) with a
small translational offset (to shift the viewpoint laterally and vertically).
At each pose, the BodySearchScanner module fuses the RealSense depth and
skeleton data into 3D keypoint estimates of the patient's torso landmarks
(shoulders, hips, sternum). The scan produces a multi-view fused torso model
from which the five canonical FAST anatomical target poses are computed.

\subsubsection{Conditional phase-three refinement}
After the ARC scan completes, the system evaluates the visibility statistics
of each keypoint. If all keypoints are detected with confidence above a
quality threshold, the FAST target poses are published immediately and the
scan is declared complete. If key keypoints---typically the shoulders or hips,
which may be occluded by patient clothing or posture---fall below the
threshold, an adaptive third scan phase is triggered. This phase adds
targeted viewpoints along the body axis (to recover occluded hips) or with
lateral offsets (to resolve asymmetric shoulder detections). The phase-three
refinement ensures robust anatomical localisation even under challenging
visibility conditions, without extending the scan duration when conditions
are favourable.

The \ac{WBC} look-at controller remains active throughout the scan, providing
the orientation component for each scan pose while the body scanner sets the
translational target. This division of labour---perception sets position,
\ac{WBC} sets orientation---ensures that the camera is always pointed at the
region of interest regardless of the arm configuration.

\subsubsection{Z1 FSM integration}
Once the scan is complete and the FAST targets are published in
$\mathcal{F}_W$, a ready signal is broadcast. The Z1 \ac{FSM}, which normally
executes its own body scanning routine, receives this signal and skips its
internal scan phase, transitioning directly to the scanning sequence. This
integration eliminates redundant scanning and ensures that the body scan
performed during approach is directly leveraged for the subsequent
ultrasound acquisition.

\subsection{Perception-Quality-Aware Velocity Control}
\label{sec:quality_velocity}

During the approach, the quality of the patient detection can vary due to
partial occlusions, changes in body posture, or motion of the platform. A
fixed approach velocity would be either overly conservative (taking
unnecessary time when detection is certain) or unsafe (charging toward an
uncertain target). \ac{TERESA} addresses this trade-off through a
QualityMonitor that continuously maps detection confidence to approach
velocity.

\subsubsection{Quality monitor}
The QualityMonitor maintains a target position in $\mathcal{F}_W$ and an
associated quality value $q \in [q_{\min},\, q_{\max}]$ measured in metres.
The target is initialised from the confidence-locked sample mean obtained
during the active search phase. It is subsequently updated only if a new
detection arrives with confidence that exceeds the current best confidence
by a minimum margin, preventing the target from drifting toward lower-quality
measurements. The quality value is computed from the posture confidence
$c$ as:
\begin{equation}
  q = q_{\max}\,(1 - c),
  \label{eq:quality_confidence}
\end{equation}
where $q_{\max}$ is a ceiling parameter. When confidence is high
($c \to 1$), the quality is low (small uncertainty radius); when confidence
drops, the quality grows, reflecting increasing positional uncertainty.
If no new confidence data arrive---for instance, during a temporary
occlusion---the quality grows linearly with time, modelling the accumulation
of uncertainty in the absence of new observations:
\begin{equation}
  q(t + \Delta t) = q(t) + \dot{q}_{\mathrm{growth}}\,\Delta t.
  \label{eq:quality_growth}
\end{equation}

\subsubsection{Velocity scaling}
The base approach velocity is scaled by a smooth function of the current
quality:
\begin{equation}
  v_{\mathrm{scale}}(q) =
  v_{\min} + \frac{1 - v_{\min}}{1 + q / q_{\mathrm{ref}}},
  \label{eq:velocity_scale}
\end{equation}
where $v_{\min} > 0$ is a floor that ensures the robot never stops
completely (avoiding deadlocks when quality is high), and $q_{\mathrm{ref}}$
is a reference quality at which the velocity is reduced to
$(1 + v_{\min})/2$. The scaling is applied multiplicatively to both the
forward velocity $v_x$ and the yaw rate $\omega_z$:
\begin{equation}
  v_x^{\mathrm{cmd}} = v_{\mathrm{scale}}(q)\,v_x, \qquad
  \omega_z^{\mathrm{cmd}} = v_{\mathrm{scale}}(q)\,\omega_z.
  \label{eq:cmd_vel_scaled}
\end{equation}

When detection is reliable ($q \approx q_{\min}$, i.e., high confidence),
$v_{\mathrm{scale}} \approx 1$ and the robot approaches at full speed.
As confidence degrades ($q$ grows), the velocity scales down smoothly,
giving the perception pipeline more time to re-acquire the target. When
no new detections arrive, the linear growth of $q$ in~\eqref{eq:quality_growth}
produces a gradual deceleration, ensuring that the robot does not charge
blindly toward a lost target.

The QualityMonitor thus closes the perception-action loop: the robot's
motion is continuously modulated by the quality of its own perceptual input,
embodying the principle that action should be proportional to certainty.
